\section{Theoretische Grundlagen}\label{sec:grundlagen}

\subsection{Sudoku als Entscheidungsproblem}\label{sec:csp}
Ein klassisches Sudoku besteht aus einem Gitter mit $9 \times 9$ Zellen, das in neun Zeilen, neun Spalten und neun Blöcke der Größe $3 \times 3$ unterteilt ist. Einige Zellen sind bereits zu Beginn mit Ziffern belegt und werden als Vorgaben bezeichnet. Das Sudoku gilt als gelöst, wenn jede Zelle eine Ziffer von 1 bis 9 enthält und jede Ziffer pro Zeile, Spalte und Block genau einmal vorkommt.

Aus Sicht der theoretischen Informatik lässt sich diese Aufgabe auf zwei Arten formulieren. Als Suchproblem wird eine konkrete Lösung verlangt. 
Als Entscheidungsproblem wird lediglich gefragt, ob zu einem vorgegebenen, teilweise gefüllten Sudoku eine gültige Lösung existiert, die Antwort lautet also entweder Ja oder Nein. Die Komplexitätstheorie betrachtet in der Regel Entscheidungsprobleme, weshalb diese Formulierung die Grundlage für Abschnitt~\ref{sec:komplexitaet} bildet. Beide Sichtweisen hängen eng zusammen, da die in dieser Arbeit betrachteten Verfahren die Frage nach der Lösbarkeit beantworten, indem sie versuchen eine Lösung zu konstruieren.

Formal werden Entscheidungsprobleme als Sprachen beschrieben. Eine Sprache ist dabei eine Menge von Wörtern, die aus den Zeichen eines festgelegten Alphabets gebildet werden \vgl{sipser2013}. Ein Alphabet bezeichnet eine endliche Menge von benötigten Zeichen, bei einem Sudoku sind die Zeichen die Ziffern von 0 bis 9. Ein Sudoku lässt sich als Wort der Länge 81 darstellen, indem die Zellen zeilenweise aneinandergereiht werden und die Ziffer 0 eine leere Zelle kennzeichnet. Die Menge aller lösbaren Sudokus bildet die Sprache.

Die Frage, ob ein gegebenes Sudoku lösbar ist, entspricht damit dem Wortproblem, also der Frage, ob das Wort zu der Sprache gehört.

Ein übliches $9 \times 9$ Sudoku kann genau eine Lösung haben, wenn 17 oder mehr Vorgaben gegeben sind \vgl{mcguire2014}. Sudokus mit wenigen oder ungünstig verteilten Vorgaben können dagegen mehrere Lösungen zulassen, widersprüchliche Vorgaben keine einzige. Für das Entscheidungsproblem ist nur relevant, ob mindestens eine zulässige Lösung existiert.

\subsection{Komplexität, Lösbarkeit und Reduktionen}\label{sec:komplexitaet}
Nicht jedes Entscheidungsproblem lässt sich algorithmisch beantworten. Turing zeigte mithilfe des Halteproblems, dass es Fragen gibt, für die kein Algorithmus existiert, der für jede Eingabe nach endlich vielen Schritten die richtige Antwort liefert \vgl{sipser2013}. Solche Probleme heißen unentscheidbar. Das Sudoku-Problem aus Abschnitt~\ref{sec:csp} gehört nicht dazu, da der Suchraum eines Sudokus endlich ist. Bei $k$ leeren Zellen existieren höchstens $9^k$ Belegungen, die sich der Reihe nach überprüfen lassen. Ein solches Verfahren findet entweder eine Lösung oder stellt fest, dass hierbei keine Lösung existiert. Sudoku ist somit also entscheidbar.

Die Komplexitätstheorie ordnet Entscheidungsprobleme nach diesem Aufwand in Klassen ein, gemessen in Abhängigkeit von der Problemgröße $n$. Die Klasse P enthält alle Probleme, die ein deterministischer Algorithmus in polynomieller Zeit löst. Die Klasse NP enthält alle Probleme, bei denen sich eine vorgeschlagene Lösung in polynomieller Zeit überprüfen lässt \vgl{sipser2013}. Gleichwertig lässt sich NP über nichtdeterministische Algorithmen beschreiben, die eine Lösung raten und anschließend in polynomieller Zeit prüfen. Jedes Problem aus P liegt auch in NP. Sudoku liegt in NP, da für ein vollständig gefülltes Sudoku alle Zeilen, Spalten und Blöcke auf doppelte Ziffern nur geprüft werden müssen.

Um Probleme hinsichtlich ihrer Schwierigkeit zu vergleichen, dient die polynomielle Reduktion. Ein Problem $A$ ist auf ein Problem $B$ reduzierbar, geschrieben $A \leq_p B$, wenn sich jede Eingabe von $A$ in polynomieller Zeit so in eine Eingabe von $B$ umformen lässt, dass beide dieselbe Antwort besitzen \vgl{sipser2013}. Ein Verfahren für $B$ löst damit auch $A$. Ein Problem heißt NP-vollständig, wenn es in NP liegt und sich jedes Problem aus NP darauf reduzieren lässt. Cook wies als erstes Problem das Erfüllbarkeitsproblem der Aussagenlogik als NP-vollständig nach \vgl{cook1971}. Für NP-vollständige Probleme ist kein Algorithmus mit polynomieller Laufzeit bekannt, ob ein solcher existiert, ist als P-NP-Problem eine der bekanntesten offenen Fragen der Informatik und eines der Millenial-Probleme. Die bekannten exakten Verfahren benötigen im ungünstigsten Fall exponentielle Laufzeit, weshalb solche Probleme bei einer großen Menge an Eingaben als praktisch nicht handhabbar gelten.

Für Sudoku gilt diese Einordnung in verallgemeinerter Form. Auf Sudokus der Größe $n^2 \times n^2$ ist die Frage nach der Lösbarkeit NP-vollständig \vgl{yato2003}. Das klassische Sudoku mit $n = 3$ besitzt eine feste Größe und ist damit streng genommen in konstanter Zeit lösbar. Die Komplexität des allgemeinen Problems erklärt dennoch, warum der Aufwand einfacher Suchverfahren bei ungünstigen Rätseln stark ansteigt.

Da Sudoku in NP liegt und SAT NP-vollständig ist, existiert eine polynomielle Reduktion von Sudoku auf SAT. Für ein Sudoku mit Kantenlänge $m$ genügen dabei $m^3$ Variablen und $O(m^4)$ Klauseln. Diese Reduktion bildet die Grundlage des zweiten Lösungsansatzes dieser Arbeit und wird in Abschnitt~\ref{sec:impl-sat} konkret umgesetzt.



\subsection{Backtracking und Heuristiken}\label{sec:backtracking}


\subsection{Aussagenlogik und das SAT-Problem}\label{sec:sat}


\subsection{Der DPLL-Algorithmus}\label{sec:dpll}
